\section{Aturan Inferensi}

\logic{Aturan inferensi} adalah pola valid untuk menarik kesimpulan dari premis \cite{rosen2019,gentzen1934}. Dalam logika, sebuah argumen dikatakan valid jika kebenaran premis-premisnya menjamin kebenaran kesimpulan; artinya, tidak mungkin premis benar sementara kesimpulan salah. Aturan inferensi menyediakan skema yang teruji untuk membangun argumen valid tanpa harus memeriksa tabel kebenaran secara eksplisit, sehingga sangat efisien untuk penalaran bertahap \cite{forallx,iep_natded}. Konsep ini dipelopori oleh Gerhard Gentzen dalam karyanya tentang natural deduction pada tahun 1934, yang memformalkan bagaimana manusia secara alami menarik kesimpulan dari premis \cite{gentzen1934,wiki_natded}. Gambar~\ref{fig:inference-flow-bab6} mengilustrasikan alur premis menuju kesimpulan untuk Modus Ponens dan Modus Tollens.

\begin{figure}[htbp]
\centering
\begin{tikzpicture}[node distance=1cm, font=\small, every node/.style={rectangle, draw=black, rounded corners=2pt, minimum height=0.7cm, align=center}]
  % Modus Ponens
  \node[fill=blue!10, minimum width=2.2cm] (mp1) at (0,1.5) {$p \rightarrow q$};
  \node[fill=blue!10, minimum width=2.2cm] (mp2) at (0,0) {$p$};
  \node[fill=green!15, minimum width=2.2cm] (mp3) at (3.2,0.75) {$q$};
  \draw[thick, ->, >=stealth] (mp1.east) -- ++(0.4,0) |- (mp3.west);
  \draw[thick, ->, >=stealth] (mp2.east) -- ++(0.4,0) |- (mp3.west);
  \node[below=0.1cm of mp2] {\textbf{Modus Ponens}};
  % Modus Tollens
  \node[fill=blue!10, minimum width=2.2cm] (mt1) at (5.5,1.5) {$p \rightarrow q$};
  \node[fill=blue!10, minimum width=2.2cm] (mt2) at (5.5,0) {$\neg q$};
  \node[fill=green!15, minimum width=2.2cm] (mt3) at (8.7,0.75) {$\neg p$};
  \draw[thick, ->, >=stealth] (mt1.east) -- ++(0.4,0) |- (mt3.west);
  \draw[thick, ->, >=stealth] (mt2.east) -- ++(0.4,0) |- (mt3.west);
  \node[below=0.1cm of mt2] {\textbf{Modus Tollens}};
\end{tikzpicture}
\caption{Alur premis menuju kesimpulan: Modus Ponens (kiri) dan Modus Tollens (kanan)}
\label{fig:inference-flow-bab6}
\end{figure}

\textbf{Modus Ponens} (dari bahasa Latin ``cara mengafirmasi'') adalah aturan inferensi paling dasar untuk implikasi \cite{rosen2019,forallx}. Bentuknya: jika kita memiliki premis $p \rightarrow q$ dan premis $p$, maka kita dapat menyimpulkan $q$. Secara semantik, jika implikasi benar dan anteseden (bagian ``jika'') benar, maka konsekuen (bagian ``maka'') harus benar. Contoh sehari-hari: ``Jika hujan, tanah basah'' dan ``Sekarang hujan'' mengarah ke kesimpulan ``Tanah basah''; dalam matematika, ``Jika $n$ genap maka $n^2$ genap'' bersama ``$n$ genap'' menghasilkan ``$n^2$ genap'' \cite{iep_natded,copi2014}.

\textbf{Modus Tollens} (dari bahasa Latin ``cara menolak'') bekerja dengan menolak konsekuen untuk menolak anteseden \cite{rosen2019,copi2014}. Bentuknya: jika $p \rightarrow q$ benar dan $\neg q$ benar, maka kita menyimpulkan $\neg p$. Berbeda dengan modus ponens yang mengafirmasi anteseden, modus tollens mengafirmasi negasi konsekuen sehingga memungkinkan kita menyimpulkan negasi anteseden. Contoh: ``Jika hujan, tanah basah'' dan ``Tanah tidak basah'' mengarah ke ``Tidak hujan''; kesalahan umum adalah memakai ``$\neg p$'' sebagai premis untuk menyimpulkan ``$\neg q$'', yang tidak valid (fallacy of denying the antecedent) \cite{levin_discrete,rosen2019}.

\textbf{Silogisme Hipotesis} memungkinkan kita merantai implikasi \cite{rosen2019,levin_discrete}. Dari $p \rightarrow q$ dan $q \rightarrow r$ kita menyimpulkan $p \rightarrow r$. Aturan ini berguna ketika kita memiliki serangkaian implikasi berurutan dan ingin menyimpulkan implikasi dari titik awal ke titik akhir. Contoh: ``Jika saya bangun pagi, saya sarapan'' dan ``Jika saya sarapan, saya berkonsentrasi baik'' menghasilkan ``Jika saya bangun pagi, saya berkonsentrasi baik''. Dalam pembuktian matematika, silogisme hipotesis sering dipakai secara implisit ketika kita membangun rantai langkah $A \Rightarrow B \Rightarrow C \Rightarrow D$ \cite{forallx,copi2014}.

\textbf{Silogisme Disjungtif} mengandalkan disjungsi dan negasi salah satu disjunct \cite{rosen2019}. Dari $p \vee q$ dan $\neg p$ kita menyimpulkan $q$; demikian pula dari $p \vee q$ dan $\neg q$ kita menyimpulkan $p$. Ide dasarnya: jika salah satu dari $p$ atau $q$ pasti benar, dan kita tahu $p$ salah, maka $q$ harus benar. Contoh: ``Dia di rumah atau di kantor'' dan ``Dia tidak di rumah'' memberi kesimpulan ``Dia di kantor''. Aturan ini berguna dalam analisis kasus dan reasoning by cases \cite{levin_discrete,cmu_natded}. Tabel~\ref{tab:aturan-inferensi} merangkum aturan inferensi utama beserta premis dan kesimpulan \cite{rosen2019,levin_discrete,brilliant_rules_inference,wiki_rules_of_inference}.

\begin{table}[htbp]
\centering
\small
\begin{tabular}{|>{\raggedright\arraybackslash}p{2.5cm}|>{\raggedright\arraybackslash}p{3.2cm}|>{\raggedright\arraybackslash}p{2.35cm}|>{\raggedright\arraybackslash}p{3.2cm}|}
\hline
\textbf{Aturan} & \textbf{Premis} & \textbf{Kesimpulan} & \textbf{Kapan dipakai} \\
\hline
Modus Ponens & $p \rightarrow q$, $p$ & $q$ & Mengafirmasi anteseden \\
\hline
Modus Tollens & $p \rightarrow q$, $\neg q$ & $\neg p$ & Menolak konsekuen $\Rightarrow$ menolak anteseden \\
\hline
Silogisme Hipotesis & $p \rightarrow q$, $q \rightarrow r$ & $p \rightarrow r$ & Merantai implikasi \\
\hline
Silogisme Disjungtif & $p \vee q$, $\neg p$ & $q$ & Salah satu disjunct salah $\Rightarrow$ yang lain benar \\
\hline
Addition & $p$ & $p \vee q$ & Menambah disjungsi \\
\hline
Simplification & $p \wedge q$ & $p$ (atau $q$) & Mengambil salah satu konjungsi \\
\hline
Conjunction & $p$, $q$ & $p \wedge q$ & Menggabungkan dua pernyataan \\
\hline
Resolution & $p \vee q$, $\neg p \vee r$ & $q \vee r$ & Automated reasoning, SAT \\
\hline
\end{tabular}
\caption{Ringkasan aturan inferensi: premis dan kesimpulan}
\label{tab:aturan-inferensi}
\end{table}

Selain keempat aturan utama di atas, terdapat aturan tambahan yang sering digunakan \cite{rosen2019,cmu_natded}. \textbf{Addition}: dari $p$ kita dapat menyimpulkan $p \vee q$ (menambah disjungsi). \textbf{Simplification}: dari $p \wedge q$ kita menyimpulkan $p$ (atau $q$). \textbf{Conjunction}: dari $p$ dan $q$ kita menyimpulkan $p \wedge q$. \textbf{Resolution}: dari $p \vee q$ dan $\neg p \vee r$ kita menyimpulkan $q \vee r$; aturan ini menjadi dasar banyak algoritma automated theorem proving dan SAT solving \cite{wikipedia_sat,cmu_sat}. Menguasai kombinasi aturan-aturan ini memungkinkan kita membangun argumen kompleks dari langkah-langkah elementer yang valid.

\begin{contoh}
Argumen berikut menggunakan modus ponens dan silogisme hipotesis. Premis: (1) ``Jika mahasiswa lulus ujian, ia mendapat nilai A'' ($p \rightarrow q$), (2) ``Jika ia mendapat nilai A, ia diterima beasiswa'' ($q \rightarrow r$), (3) ``Mahasiswa itu lulus ujian'' ($p$). Dari (1) dan (3) dengan modus ponens: $q$ (``ia mendapat nilai A''). Dari $q \rightarrow r$ dan $q$ dengan modus ponens: $r$ (``ia diterima beasiswa''). Alternatif: dari (1) dan (2) dengan silogisme hipotesis diperoleh $p \rightarrow r$, lalu dengan (3) dan modus ponens diperoleh $r$.
\end{contoh}

\begin{contoh}
\textbf{Simulasi Modus Ponens dan Modus Tollens dengan Python} \cite{python_docs_boolean,rosen2019}

Fungsi berikut menerima nilai kebenaran premis dan menampilkan kesimpulan sesuai aturan. Untuk Modus Ponens: jika $p \to q$ benar dan $p$ benar, maka kesimpulan $q$; untuk Modus Tollens: jika $p \to q$ benar dan $\neg q$ benar, maka kesimpulan $\neg p$.

\begin{lstlisting}[style=pythonstyle, caption={Pengecekan kesimpulan Modus Ponens dan Modus Tollens}]
def implikasi(p, q):
    """p -> q (False hanya jika p True dan q False)."""
    return (not p) or q

def modus_ponens(p_impl_q, p):
    """Jika p->q dan p benar, kesimpulan q."""
    if p_impl_q and p:
        return True   # q harus True
    return None       # tidak dapat menyimpulkan

def modus_tollens(p_impl_q, not_q):
    """Jika p->q dan not q benar, kesimpulan not p."""
    if p_impl_q and not_q:
        return True   # not p benar
    return None

# Contoh: "Jika hujan, tanah basah". p=hujan, q=tanah basah
p, q = True, True   # hujan dan tanah basah
print("Modus Ponens:", modus_ponens(implikasi(p, q), p))  # True -> q
p, q = False, False # tidak hujan, tanah tidak basah
print("Modus Tollens:", modus_tollens(True, True))  # not q -> not p
\end{lstlisting}
\end{contoh}

\begin{konsep}
Ringkasan aturan inferensi utama: Modus Ponens ($p \rightarrow q$, $p$ $\therefore$ $q$), Modus Tollens ($p \rightarrow q$, $\neg q$ $\therefore$ $\neg p$), Silogisme Hipotesis ($p \rightarrow q$, $q \rightarrow r$ $\therefore$ $p \rightarrow r$), Silogisme Disjungtif ($p \vee q$, $\neg p$ $\therefore$ $q$). Semua aturan ini valid: jika premis benar, kesimpulan harus benar. Jangan keliru dengan fallacy seperti affirming the consequent ($p \rightarrow q$, $q$ $\therefore$ $p$) atau denying the antecedent ($p \rightarrow q$, $\neg p$ $\therefore$ $\neg q$), yang tidak valid \cite{mit_proof_methods,brilliant_rules_inference}.
\end{konsep}

\begin{itemize}
  \item \textbf{Modus Ponens}: Jika $p \rightarrow q$ dan $p$, maka $q$.
  \item \textbf{Modus Tollens}: Jika $p \rightarrow q$ dan $\neg q$, maka $\neg p$.
  \item \textbf{Silogisme Hipotesis}: Jika $p \rightarrow q$ dan $q \rightarrow r$, maka $p \rightarrow r$.
  \item \textbf{Silogisme Disjungtif}: Jika $p \vee q$ dan $\neg p$, maka $q$.
\end{itemize}
